\begin{abstract}
Evaluating the security posture of Large Language Model (LLM) deployment stacks is a critical challenge in modern artificial intelligence security. Traditional vulnerability management frameworks---such as the Common Vulnerability Scoring System (CVSS) and component-level checklists---critically assume that software components can be evaluated in isolation. In real-world agentic and RAG-based LLM ecosystems, this assumption is systematically violated as attackers exploit complex topologies, chaining seemingly low-risk vulnerabilities (e.g., indirect prompt injection) with downstream tools (e.g., SQL execution) to achieve catastrophic compromises. Na\"ively applying independent scoring methods to deeply integrated LLM stacks yields inflated risk assessments, misaligned mitigation priorities, and a failure to capture compositional attack paths.

We propose \textbf{Compositional Attack Path Scoring (CAPS)}, a novel security framework specifically engineered to quantify end-to-end multi-hop risks in LLM architectures. CAPS uniquely integrates three synergistic capabilities: (i) directed graph topological modeling, which maps the full deployment stack from attacker entry points to high-value internal assets; (ii) dynamic mitigation attenuation, which calculates the ``Effective Exploitability'' of individual nodes based on the presence of targeted security guardrails; and (iii) a mathematical compositional path engine that optimizes risk scoring via an explicit exponential decay factor, reflecting the real-world friction of traversing complex trust boundaries. Furthermore, CAPS provides an automated Return on Investment (ROI) recommendation engine to rank mitigation strategies by their systemic risk reduction.

Empirical evaluation on a suite of standardized LLM architectures (including RAG Chatbots, Autonomous Coding Agents, and Enterprise Model Routers) demonstrates that CAPS significantly improves risk calibration. Against the Autonomous Agent benchmark, CAPS calculates a realistic critical path score of 51.3, successfully correcting the naive 85.0 overestimation produced by traditional component-level CVSS scoring. CAPS establishes a new, rigorous benchmark for quantitative vulnerability management in complex, agentic LLM environments.
\end{abstract}

\begin{keywords}
Large Language Models, AI Security, Attack Path Analysis, Threat Modeling, Compositional Scoring, Vulnerability Management
\end{keywords}
